%=====================================================================
\chapter{Representations and Search in Python}\label{ch:python}
%=====================================================================
\locator{1}{Part I --- Foundations of Intelligent Systems}

\begin{outcomes}
By the end of this chapter you will be able to:
\begin{tight}
  \item \textbf{Represent} a state as an immutable, hashable value, and
        \textbf{explain} why every search algorithm in this book requires it.
  \item \textbf{Implement} the \texttt{Problem} interface --- initial state, goal
        test, actions, result, step cost --- that every later chapter subclasses.
  \item \textbf{Choose} the right frontier structure for a search strategy, and
        \textbf{justify} the choice by its cost: \texttt{deque} for first-in
        first-out, list for last-in first-out, \texttt{heapq} for best-first.
  \item \textbf{Encode} a project's task dependencies as a graph and
        \textbf{separate} the problem's data from the algorithm that searches it.
  \item \textbf{Use} generators for successor functions and \texttt{lru\_cache}
        for repeated subproblems, and \textbf{measure} what each one saves.
\end{tight}
\end{outcomes}

\begin{prereq}
\chref{agents} (agent, percept, action, goal-based agent). Object Oriented
Programming: classes, inheritance, abstract methods, and the difference between a
value and a reference. Appendix~A is the reference for every standard-library call
used here.
\end{prereq}

\begin{keyterms}
state $\bullet$ immutability $\bullet$ hashability $\bullet$ \texttt{frozenset}
$\bullet$ \texttt{namedtuple} $\bullet$ the \texttt{Problem} interface $\bullet$
successor function $\bullet$ node $\bullet$ frontier $\bullet$ explored set
$\bullet$ \texttt{deque} $\bullet$ \texttt{heapq} $\bullet$ priority queue
$\bullet$ adjacency dictionary $\bullet$ generator $\bullet$ memoisation $\bullet$
\texttt{lru\_cache}
\end{keyterms}

\section{Why This Chapter Exists}

The prerequisite for this course is Object Oriented Programming, so you can already
write Python. This chapter is not a Python tutorial. It is the chapter that fixes
the four data structures the rest of the book is written in, so that Chapters~4
to~9 can say ``push the successor onto the frontier'' and mean something precise.

Every search algorithm in Part~II is about twenty lines long. They are short only
because the representation carries the weight. Get the representation wrong and the
algorithm triples in size and starts to produce wrong answers for reasons that are
extremely hard to see. That is the whole content of this chapter, and it is worth a
week.

\begin{conceptbox}
Here is the judgement to take away, before the details. \emph{A state is a value,
not an object.} Two states that describe the same situation must be equal and must
hash the same, whether or not they are the same object in memory. If you find
yourself writing \texttt{if state1.tasks == state2.tasks and state1.hours ==
state2.hours}, you have already lost: the language will do that for you, correctly
and once, if the state is a tuple or a \texttt{frozenset}.
\end{conceptbox}

\section{A State Is a Value}

Consider the state of a release: which components have been coded, reviewed,
tested and deployed. The obvious first attempt is a class with mutable attributes.
It is wrong for three reasons, and all three matter.

\begin{alertbox}[Three reasons a mutable state object breaks search]
\begin{tight}
  \item \textbf{You cannot put it in a set.} Search must ask ``have I seen this
        state before?'' millions of times. That question is a set membership test,
        and Python can only put hashable objects in a set. A class with mutable
        attributes is hashable only by identity, so two identical situations look
        different and the search explores the same state again and again ---
        usually forever.
  \item \textbf{Mutation corrupts the frontier.} If a successor is produced by
        changing the parent in place, every node already sitting in the frontier
        that shares that object silently changes too. The bug appears as a search
        that returns a path that does not work when executed.
  \item \textbf{You cannot reconstruct the path.} Solutions are recovered by
        walking parent pointers back to the root. If the states were mutated on the
        way down, the ancestors no longer hold the values they had.
\end{tight}
\end{alertbox}

So states are built from Python's immutable, hashable types:
\texttt{tuple}, \texttt{frozenset}, \texttt{str}, numbers, and
\texttt{namedtuple}. \dsref{state-repr} shows the three representations this book
uses and when each is right.

\dsfig{%
\begin{tikzpicture}[font=\scriptsize,
  card/.style={rounded corners=3pt, draw=#1, line width=1.1pt, fill=#1!7,
               text width=4.3cm, align=flush left, inner sep=6pt,
               text=#1!50!black, minimum height=4.0cm},
  hd/.style={rounded corners=2pt, fill=#1, text=white, font=\scriptsize\bfseries,
             inner sep=3pt, text width=4.1cm, align=center},
  cc/.style={font=\ttfamily\scriptsize, text=neutral}]

\node[card=primary] (c1) at (0,0)
  {\\[2pt]\textbf{Use when:} the state is a \emph{set of things that have
   happened}, and order is meaningless.\\[4pt]
   \texttt{frozenset(\{"code A",}\\
   \texttt{~~"code B", "review A"\})}\\[4pt]
   \textbf{Why:} set equality ignores order, so two paths that did the same work
   in a different order collapse to one state --- a large saving, for free.\\[4pt]
   \emph{Chapters 4, 5, 9}};

\node[card=secondary, right=0.55cm of c1] (c2)
  {\\[2pt]\textbf{Use when:} the state has \emph{fixed named slots}, each with a
   value.\\[4pt]
   \texttt{Node(task=3,}\\
   \texttt{~~hours=12, owner="AY")}\\[4pt]
   \textbf{Why:} a \texttt{namedtuple} is a tuple, so it is hashable and
   comparable, but the fields are readable by name. Printing one is
   self-documenting.\\[4pt]
   \emph{Chapters 2, 7, 14}};

\node[card=greenhl, right=0.55cm of c2] (c3)
  {\\[2pt]\textbf{Use when:} the state is an \emph{assignment} of values to a
   known list of variables.\\[4pt]
   \texttt{("S1", None, "S3", None)}\\[4pt]
   \textbf{Why:} position carries the variable, so a partial assignment is just a
   tuple with \texttt{None} in the unassigned slots. Cheap to copy with one slot
   changed.\\[4pt]
   \emph{Chapters 7, 8, 11}};

\node[hd=primary]   at (c1.north) {FROZENSET};
\node[hd=secondary] at (c2.north) {NAMEDTUPLE};
\node[hd=greenhl]   at (c3.north) {TUPLE OF SLOTS};

\node[rounded corners=3pt, fill=accent!8, draw=accent, line width=0.9pt,
      text=accent!55!black, font=\scriptsize, text width=13.6cm,
      align=flush left, inner sep=6pt] at (4.85,-2.85)
  {\textbf{The one test that catches every mistake here.} Build the same situation
   twice, by two different routes, and assert that the two states are equal and
   that a two-element set of them has length one. If that assertion fails, no
   search built on the representation will terminate reliably --- and the failure
   will look like an algorithm bug for days.};
\end{tikzpicture}}%
{Three state representations, and the situation each one is for. All three are
immutable and hashable; that is the property the algorithms depend on.}
{state-repr}

\section{The Problem Interface}

Every search problem in Chapters~4 to~9 is the same five things. Fixing them as one
small class means the algorithms can be written once and reused unchanged.

\begin{definitionbox}[A search problem]
A search problem is a five-tuple:
\begin{tight}
  \item an \term{initial state};
  \item a \term{goal test}, a predicate on states --- not necessarily a single
        goal state, because ``every component deployed'' describes many states;
  \item an \term{actions} function giving the actions legal in a state;
  \item a \term{result} function giving the state that follows from applying an
        action --- the \term{successor function}, and the only place the rules of
        the domain live;
  \item a \term{step cost} for taking an action in a state.
\end{tight}
Nothing else. Notably \emph{not} part of the problem: how to search it. That
separation is the point.
\end{definitionbox}

\begin{examplebox}[The interface, once, for the whole book]
\begin{lstlisting}[language=Python]
class Problem:
    """What every search algorithm in Chapters 4-9 is given."""

    def __init__(self, initial, goal=None):
        self.initial, self.goal = initial, goal

    def actions(self, state):        # -> iterable of actions
        raise NotImplementedError

    def result(self, state, action):  # -> new state (never mutated)
        raise NotImplementedError

    def goal_test(self, state):
        return state == self.goal

    def step_cost(self, state, action, nxt):
        return 1                      # override for weighted problems
\end{lstlisting}
A concrete problem overrides three methods and inherits the rest. Chapter~4's
task-ordering problem, Chapter~5's sprint problem and Chapter~9's release planner
are all subclasses of this, and the identical \texttt{breadth\_first} function
solves all three.
\end{examplebox}

A \term{node} is not a state. A node is a state plus the bookkeeping the search
needs: the parent it came from, the action that produced it, and the cost so far.
States get compared and stored in sets; nodes get put on frontiers and walked
backwards to build the solution. Confusing the two is the second most common bug in
this material, after mutable states.

\section{Frontiers: Three Structures, Three Strategies}

The \term{frontier} holds the nodes that have been generated but not yet expanded.
Remarkably, the choice of frontier structure \emph{is} the choice of search
strategy: the algorithm body does not change at all between breadth-first,
depth-first and best-first search. Only the container does. \dsref{frontiers} shows
the three, with the operation costs that decide the matter.

\dsfig{%
\begin{tikzpicture}[font=\scriptsize,
  fr/.style={rounded corners=3pt, draw=#1, line width=1.1pt, fill=#1!7,
             text width=4.35cm, align=flush left, inner sep=6pt,
             text=#1!50!black, minimum height=3.5cm},
  hd/.style={rounded corners=2pt, fill=#1, text=white, font=\scriptsize\bfseries,
             inner sep=3pt, text width=4.15cm, align=center},
  cell/.style={draw=#1!60, fill=#1!12, minimum width=0.42cm,
               minimum height=0.42cm, font=\tiny, text=#1!50!black},
  ar/.style={-{Stealth[length=1.6mm]}, #1, line width=0.8pt}]

\node[fr=primary] (f1) at (0,0)
  {\\[3pt]\texttt{collections.deque}\\[3pt]
   \texttt{popleft()} takes the \emph{oldest} node.\\[4pt]
   Append and pop from either end in constant time.\\[4pt]
   \textbf{Never use a plain list} for this: \texttt{list.pop(0)} is linear, and a
   breadth-first search on a list is quadratic for no reason.\\[3pt]
   \emph{Breadth-first, Chapter 4}};

\node[fr=secondary, right=0.5cm of f1] (f2)
  {\\[3pt]\texttt{list}\\[3pt]
   \texttt{pop()} takes the \emph{newest} node.\\[4pt]
   Append and pop from the right in constant time.\\[4pt]
   A plain list is exactly right here, and the recursive form of the same search
   uses Python's own call stack instead.\\[3pt]
   \emph{Depth-first, Chapter 4}};

\node[fr=greenhl, right=0.5cm of f2] (f3)
  {\\[3pt]\texttt{heapq} on a list\\[3pt]
   \texttt{heappop()} takes the \emph{cheapest} node.\\[4pt]
   Push and pop in logarithmic time; the heap is a plain list of
   \texttt{(priority, counter, node)} tuples.\\[4pt]
   The counter breaks ties so two equal priorities never compare the nodes.\\[3pt]
   \emph{Uniform-cost and A*, Chapters 4--5}};

\node[hd=primary]   at (f1.north) {FIFO QUEUE};
\node[hd=secondary] at (f2.north) {LIFO STACK};
\node[hd=greenhl]   at (f3.north) {PRIORITY QUEUE};

% little queue diagrams underneath
\foreach \i/\lab in {0/1,1/2,2/3,3/4}{
  \node[cell=primary] at (-1.45+\i*0.46,-2.35) {\lab};}
\draw[ar=primary] (-1.95,-2.35) -- ++(-0.45,0)
  node[left, font=\tiny, text=primary] {out};
\draw[ar=primary] (0.42,-2.35) -- ++(0.45,0)
  node[right, font=\tiny, text=primary] {in};

\foreach \i/\lab in {0/1,1/2,2/3,3/4}{
  \node[cell=secondary] at (3.4+\i*0.46,-2.35) {\lab};}
\draw[ar=secondary] (5.26,-2.05) -- ++(0.0,0.42)
  node[above, font=\tiny, text=secondary] {in/out};

\foreach \i/\lab in {0/1,1/3,2/7,3/9}{
  \node[cell=greenhl] at (8.25+\i*0.46,-2.35) {\lab};}
\draw[ar=greenhl] (7.75,-2.35) -- ++(-0.45,0)
  node[left, font=\tiny, text=greenhl!50!black] {min};

\node[font=\tiny\itshape, text=neutral, text width=13.6cm, align=flush left]
  at (4.85,-3.15)
  {The search body is identical in all three cases: pop a node, test the goal,
   push its successors. Chapter~4 writes it once and changes only this container.};
\end{tikzpicture}}%
{Three frontier structures. Choosing the container chooses the strategy, and the
operation costs are the reason each one is matched to its search.}{frontiers}

\begin{alertbox}[Two frontier mistakes that cost real time]
\textbf{Using a list as a queue.} \texttt{frontier.pop(0)} on a list moves every
remaining element one place left. It is $O(n)$ per pop, so a breadth-first search
generating $n$ nodes does $O(n^{2})$ pointless work. The cost is invisible on this
chapter's nine tasks and decisive on anything the size of a real backlog. Use
\texttt{deque.popleft()}.

\smallskip
\textbf{Pushing bare tuples onto a heap.} \texttt{heappush(h, (cost, node))}
compares the second element whenever two costs are equal, and if \texttt{node} is a
class without ordering, Python raises \texttt{TypeError} --- intermittently, only
on inputs that happen to tie. Always push
\texttt{(priority, next(counter), node)}, with a monotonically increasing counter
in the middle. The tie is then broken by insertion order and the node is never
compared.
\end{alertbox}

\section{Representing the Project}

All the domain data in this book is one small software project: nine tasks with a
duration and a set of dependencies each. It lives in plain dictionaries,
deliberately, so that every lab can be read without reference to a data file.

\begin{examplebox}[The project, as data]
\begin{lstlisting}[language=Python]
# task -> (duration in days, set of tasks that must finish first)
TASKS = {
    "spec":      (3, set()),
    "schema":    (2, {"spec"}),
    "api":       (5, {"schema"}),
    "ui":        (4, {"spec"}),
    "auth":      (3, {"schema"}),
    "tests":     (4, {"api", "auth"}),
    "docs":      (2, {"api"}),
    "uat":       (3, {"tests", "ui"}),
    "release":   (1, {"uat", "docs"}),
}

def ready(done):
    """Tasks whose dependencies are all satisfied -- the successor
    function of Chapter 4, in one line."""
    return {t for t, (_, deps) in TASKS.items()
            if t not in done and deps <= done}
\end{lstlisting}
Two things to notice. \texttt{deps <= done} is subset containment, which reads as
``every dependency is done'' --- this is why dependencies are sets. And
\texttt{ready} returns the legal actions for a state, so it \emph{is} the
\texttt{actions} method of a \texttt{Problem}; Chapter~4 does no more than wrap it.
\end{examplebox}

\section{Generators and Memoisation}

Two standard-library habits pay for themselves repeatedly in this course.

A \term{generator} successor function yields states one at a time instead of
building a list. In search this matters because most successors are never used: a
depth-first search that finds its goal in the first branch has generated one
successor of each node on the path, not all of them. Writing \texttt{yield} instead
of \texttt{result.append} is the whole change.

\term{Memoisation} caches a function's result against its arguments.
\texttt{functools.lru\_cache} does it in one decorator, and it requires exactly the
property Section~2 insisted on: the arguments must be hashable. Chapter~5's
heuristic is called many times on the same state, Chapter~8's minimax re-reaches
the same position by different move orders, and Chapter~11's model counting
revisits the same partial assignment. In all three, one decorator changes the
running time by an order of magnitude and nothing else in the code.

\begin{worked}[Why the state representation decides whether search terminates]
Two representations of ``\texttt{spec} and \texttt{schema} are done''. Work out
what each does to the search.

\smallskip
\textbf{Representation A --- a list, in completion order.}
\begin{tight}
  \item Route 1: do \texttt{spec}, then \texttt{schema} $\;\Rightarrow\;$
        \texttt{["spec", "schema"]}
  \item Route 2: not possible here, since \texttt{schema} depends on \texttt{spec}
        --- so extend the example by one task. With \texttt{ui} also done:
        \texttt{["spec", "schema", "ui"]} and \texttt{["spec", "ui", "schema"]}.
\end{tight}
These two lists are not equal. A list is also unhashable, so they cannot go in a
set at all. Both routes are explored separately, and everything below them is
explored twice.

\smallskip
\textbf{Counting the damage.} With $n$ mutually independent tasks completed in any
order, the number of distinct \emph{orderings} is $n!$ while the number of
distinct \emph{situations} is $1$. The dependencies in this project rule out most
orderings, so the real gap is smaller than $n!$ but still large: the lab counts
$184$ reachable sequences against $22$ reachable states, a factor of $8$.

\smallskip
\textbf{Representation B --- a \texttt{frozenset}.}
\[
  \texttt{frozenset(\{"spec","schema","ui"\})}
  \;=\;
  \texttt{frozenset(\{"spec","ui","schema"\})}
\]
Equal, and equal hashes. So the second route arrives at a state already in the
explored set and is cut off immediately.

\smallskip
\textbf{The measurement.} This chapter's lab reports both numbers on this project:
$22$ states reachable as sets against $184$ paths as sequences. The search is the
same search. \emph{The representation bought a factor of eight, and not one line of
algorithm changed.}

\smallskip
And the factor grows with the project. Add one task that depends on nothing and is
required by nothing, and it can be inserted at any point in every existing
sequence: the sequence count multiplies by roughly the plan length, while the state
count merely doubles. On a realistic backlog of forty loosely coupled stories the
ratio is no longer eight but astronomical, which is why this is a representation
decision and not a micro-optimisation.

\smallskip
\textbf{The lesson to carry forward.} When a search is too slow, look at the state
representation \emph{before} looking at the algorithm. Ask: does my state record
anything the goal test does not care about? Here, the order tasks were finished in
is such a thing. Removing it is called \emph{abstraction}, and it is the most
effective single optimisation in this book.
\end{worked}

\begin{pitfall}
\begin{tight}
  \item \textbf{A state class with mutable attributes.} Unhashable in practice,
        hashable by identity at worst, and the resulting search never terminates.
        Use a tuple, \texttt{namedtuple} or \texttt{frozenset}.
  \item \textbf{Mutating a state to make a successor.} \texttt{done.add(t)} changes
        the parent that other frontier nodes point at. Write
        \texttt{done | \{t\}}, which builds a new set.
  \item \textbf{Recording in the state something the goal test ignores.} Completion
        order, a timestamp, a counter of loop iterations. Each one multiplies the
        state space.
  \item \textbf{\texttt{list.pop(0)} as a queue.} Quadratic. Use
        \texttt{collections.deque}.
  \item \textbf{Heap entries without a tie-breaker.} Works until two priorities
        are equal, then raises \texttt{TypeError} on unorderable nodes ---
        intermittently, which is the worst way for a bug to behave.
  \item \textbf{\texttt{lru\_cache} on a function taking a list.} Raises
        \texttt{TypeError: unhashable type}. This is the same discipline as the
        states, enforced by the standard library.
\end{tight}
\end{pitfall}

%---------------------------------------------------------------------
\begin{lab}[The Frontier and the Graph: The Course Toolkit]
Everything later chapters import. Build it once, assert that it behaves as
documented, and measure the claim of the worked example.
\begin{lstlisting}[language=Python]
"""Chapter 3 lab -- the course toolkit.

The Problem interface, three frontiers, the project graph, and the
measurement that justifies using frozensets for states. Every later
chapter builds on exactly these shapes.
"""

import heapq
import itertools
from collections import deque, namedtuple

# --- the project, as data ------------------------------------------
TASKS = {
    "spec":    (3, set()),
    "schema":  (2, {"spec"}),
    "api":     (5, {"schema"}),
    "ui":      (4, {"spec"}),
    "auth":    (3, {"schema"}),
    "tests":   (4, {"api", "auth"}),
    "docs":    (2, {"api"}),
    "uat":     (3, {"tests", "ui"}),
    "release": (1, {"uat", "docs"}),
}
ALL = frozenset(TASKS)


# --- the interface every later chapter subclasses -------------------
class Problem:
    """A search problem: five things, and no search strategy."""

    def __init__(self, initial, goal=None):
        self.initial, self.goal = initial, goal

    def actions(self, state):
        raise NotImplementedError

    def result(self, state, action):
        raise NotImplementedError

    def goal_test(self, state):
        return state == self.goal

    def step_cost(self, state, action, nxt):
        return 1


class TaskOrder(Problem):
    """Finish every task, respecting dependencies.

    A state is the frozenset of finished tasks -- nothing else. The
    order they were finished in is deliberately not recorded.
    """

    def __init__(self):
        super().__init__(frozenset(), ALL)

    def actions(self, state):
        for t, (_, deps) in TASKS.items():
            if t not in state and deps <= state:
                yield t                      # a generator, not a list

    def result(self, state, action):
        return state | {action}              # new value, never mutated

    def step_cost(self, state, action, nxt):
        return TASKS[action][0]              # duration in days


Node = namedtuple("Node", "state parent action cost")


def path(node):
    """Walk parent pointers back to the root."""
    out = []
    while node.parent is not None:
        out.append(node.action)
        node = node.parent
    return list(reversed(out))


# --- the three frontiers, behind one small interface ----------------
class FIFO:
    def __init__(self):
        self.q = deque()

    def push(self, node):
        self.q.append(node)

    def pop(self):
        return self.q.popleft()          # O(1); list.pop(0) is O(n)

    def __len__(self):
        return len(self.q)


class LIFO:
    def __init__(self):
        self.q = []

    def push(self, node):
        self.q.append(node)

    def pop(self):
        return self.q.pop()

    def __len__(self):
        return len(self.q)


class Cheapest:
    """A heap of (priority, tie, node). The counter keeps the heap
    from ever comparing two nodes."""

    def __init__(self):
        self.h, self.tie = [], itertools.count()

    def push(self, node):
        heapq.heappush(self.h, (node.cost, next(self.tie), node))

    def pop(self):
        return heapq.heappop(self.h)[2]

    def __len__(self):
        return len(self.h)


def search(problem, Frontier):
    """One body; the container decides the strategy. Chapter 4 proves
    this by calling it three times."""
    start = Node(problem.initial, None, None, 0)
    frontier, seen = Frontier(), {problem.initial}
    frontier.push(start)
    expanded, peak = 0, 1
    while len(frontier):
        node = frontier.pop()
        if problem.goal_test(node.state):
            return path(node), node.cost, expanded, peak
        expanded += 1
        for a in problem.actions(node.state):
            nxt = problem.result(node.state, a)
            if nxt not in seen:
                seen.add(nxt)
                frontier.push(Node(nxt, node, a,
                                   node.cost + problem.step_cost(
                                       node.state, a, nxt)))
        peak = max(peak, len(frontier))
    return None, None, expanded, peak


# --- the measurement of the worked example -------------------------
def count_states():
    """Reachable situations, when a state is a set of finished tasks."""
    seen, stack = {frozenset()}, [frozenset()]
    while stack:
        s = stack.pop()
        for t, (_, deps) in TASKS.items():
            if t not in s and deps <= s:
                n = s | {t}
                if n not in seen:
                    seen.add(n)
                    stack.append(n)
    return len(seen)


def count_orderings():
    """Reachable sequences, when order is (wrongly) part of the
    state."""
    total = 0
    stack = [frozenset()]
    while stack:
        s = stack.pop()
        total += 1
        for t, (_, deps) in TASKS.items():
            if t not in s and deps <= s:
                stack.append(s | {t})
    return total


if __name__ == "__main__":
    # 1. the representation test of Figure 3.1
    a = frozenset({"spec", "schema", "ui"})
    b = frozenset({"spec", "ui", "schema"})
    assert a == b and len({a, b}) == 1
    print("states equal by two routes:", a == b)

    # 2. the three frontiers on the same problem
    p = TaskOrder()
    print()
    print("frontier      plan len   days   expanded   peak")
    print("-" * 50)
    for name, F in (("FIFO     ", FIFO), ("LIFO     ", LIFO),
                    ("Cheapest ", Cheapest)):
        plan, cost, exp, peak = search(p, F)
        print("%s %9d %6d %10d %6d" % (name, len(plan), cost, exp, peak))

    # 3. the factor the representation buys
    s, o = count_states(), count_orderings()
    print()
    print("situations (frozenset states):", s)
    print("sequences  (order in state)  :", o)
    print("factor saved                 : %.0fx" % (o / s))
    assert o > 5 * s, "sets should collapse many sequences into one"

    # 4. every plan is a legal ordering
    plan, _, _, _ = search(p, FIFO)
    done = set()
    for t in plan:
        assert TASKS[t][1] <= done, "dependency violated: " + t
        done.add(t)
    assert done == set(TASKS)
    print()
    print("plan is a legal topological order:", " -> ".join(plan))
\end{lstlisting}
All three frontiers find a legal ordering of all nine tasks, and all three report a
cost of $27$ days --- every task has to be done, so the total duration does not
depend on the order. That is precisely why cost only becomes interesting in
\chref{informed}, where not every action is taken. Where the three differ is work
done: breadth-first and uniform-cost expand $21$ states with a peak frontier of $6$
and $5$, depth-first expands $9$ with a peak of $4$. The representation
measurement comes out at $22$ states against $184$ sequences.

Now break it on purpose: change \texttt{result} to return
\texttt{state + (action,)} over a tuple state, and watch \texttt{expanded} climb
towards the sequence count while the answer stays the same. That experiment is the
chapter.
\end{lab}

%---------------------------------------------------------------------
\begin{checkpoint}
\begin{tightnum}
  \item Why can a Python \texttt{set} not hold a state represented as a list, and
        what goes wrong in a search when the explored set cannot be built?
  \item A colleague's best-first search raises a \texttt{TypeError} about
        comparing two \texttt{Node} objects, but only on about one input in
        twenty. Explain the bug and give the one-line fix.
  \item A state records the set of finished tasks \emph{and} the number of search
        steps taken so far. The goal test looks only at the finished tasks. What
        does the extra field do to the size of the state space, and why?
\end{tightnum}
\end{checkpoint}

%---------------------------------------------------------------------
\begin{chaptersummary}
\begin{tight}
  \item A state is a \emph{value}: immutable and hashable, so that two descriptions
        of the same situation are equal and hash alike. Use \texttt{frozenset} for
        sets of completed things, \texttt{namedtuple} for fixed named slots, and a
        tuple of slots for assignments.
  \item A \term{node} is a state plus parent, action and cost. States go in sets;
        nodes go on frontiers and are walked backwards to recover the solution.
  \item Every search problem is five things: initial state, goal test, actions,
        result and step cost. The strategy is not one of them --- which is why one
        search body serves every chapter of Part~II.
  \item The frontier structure \emph{is} the strategy: \texttt{deque} for
        breadth-first, list for depth-first, \texttt{heapq} for best-first. Never
        use \texttt{list.pop(0)}, and always give heap entries a tie-breaking
        counter.
  \item Successor functions are generators, because most successors are never
        needed. \texttt{lru\_cache} removes repeated subproblems for one line, and
        requires the same hashability the states already have.
  \item When a search is too slow, look first at what the state records that the
        goal test does not care about. Removing completion order collapsed $184$
        sequences into $22$ states on this chapter's project, and the saving grows
        with the size of the backlog.
\end{tight}
\end{chaptersummary}

\begin{reviewq}
  \item \textbf{(MCQ)} Which type is \emph{not} usable as a search state?
        (a)~\texttt{tuple} (b)~\texttt{frozenset} (c)~\texttt{list}
        (d)~\texttt{namedtuple}.
  \item \textbf{(MCQ)} Popping the oldest node from the frontier gives:
        (a)~depth-first search (b)~breadth-first search (c)~uniform-cost search
        (d)~best-first search.
  \item \textbf{(MCQ)} \texttt{deps <= done} where both are sets tests that:
        (a)~\texttt{done} has fewer elements (b)~every dependency is finished
        (c)~the sets are equal (d)~\texttt{deps} is sorted before \texttt{done}.
  \item \textbf{(MCQ)} The middle element of a heap entry
        \texttt{(cost, counter, node)} exists to: (a)~speed up the heap
        (b)~record insertion time for reporting (c)~stop Python comparing two
        nodes when costs tie (d)~allow duplicate states.
  \item \textbf{(Short)} State the difference between a state and a node, and name
        one thing each is used for that the other cannot do.
  \item \textbf{(Short)} Explain why making a successor by mutating the parent
        state produces a plan that fails when executed, even though the search
        reports success.
  \item \textbf{(Short)} Why is a generator the right shape for a successor
        function? Name the search strategy that benefits most, and say why.
  \item \textbf{(Applied)} Add a task \texttt{perf} of duration 2, depending on
        \texttt{api} and \texttt{ui}, and required before \texttt{release}. Write
        the new \texttt{TASKS} entry and the changed \texttt{release} entry, then
        predict --- before running it --- whether the number of reachable
        \texttt{frozenset} states goes up by more or less than a factor of two,
        and justify your prediction.
\end{reviewq}
